% TeX root = ../Main.tex

\section[Question -- {\it SIFT Algorithm}]{}

\subsection{SIFT Detection}
The Scale-Invariant Feature Transform (SIFT) detects stable keypoints in scale-space by constructing a \textbf{Difference of Gaussians (DoG)} pyramid across octaves:
\[
  \mathrm{DoG}(x, y, \sigma) = G(x, y, k\sigma) - G(x, y, \sigma)
\]
A pixel is selected as a keypoint if it is a local extremum compared to its \textbf{26 neighbors} in the 3D scale-space (8 spatial, 9 in the scale above, 9 below). 
Unstable points—specifically low-contrast points and edge responses (filtered via a 2D Hessian ratio test)—are discarded. 
Finally, a \textbf{canonical orientation} is assigned based on the peak of a local gradient orientation histogram to ensure rotation invariance.

\subsection{SIFT Descriptor}
To represent the local appearance around a keypoint:
\begin{itemize}
    \item A $16 \times 16$ pixel window is extracted at the keypoint's characteristic scale.
    \item The coordinate system and gradients are rotated according to the keypoint's canonical orientation (achieving rotation invariance).
    \item The window is divided into a $4 \times 4$ grid of cells (16 cells total).
    \item An \textbf{8-bin orientation histogram} is computed for each cell.
    \item Concatenating all 16 histograms yields a single \textbf{128-dimensional vector} ($16 \times 8 = 128$).
    \item The vector is normalized to unit length ($L_2$ norm) for illumination invariance.
\end{itemize}

\subsection{SIFT Properties}
The SIFT algorithm is robust and invariant to:
\begin{itemize}
    \item \textbf{Scale}: Via multi-scale DoG extrema detection.
    \item \textbf{Rotation}: Via canonical orientation alignment.
    \item \textbf{Illuminazione}: Via $L_2$ vector normalization.
\end{itemize}